\documentclass[12pt]{article}
\usepackage[utf8]{inputenc}
\usepackage[T1]{fontenc}
\usepackage{hyperref}
\title{Pula Yi Language Resources for K-12 Curriculum Development}
\author{Compiled by Claude Sonnet 5.5 (claude-sonnet-5-5) for Andr\'as Kornai}
\date{}
\begin{document}
\maketitle

\section{Introduction}
Pula Yi (ISO 639-3 ypl, Glottocode pula1268) is a retired language code. According to the search results, it was retired in 2008 and split into Phola (ypg), Phala (ypa) and Alo Phola (ypo), because of low intelligibility between Phola and Phala. The ``family: Bookkeeping'' field in the brief is not a language family. Glottolog labels retired codes ``Bookkeeping''. Older sources describe the Pula as a Yunnan group classified as Yi, with about 20,000 people in Honghe, Yuanjiang and nearby counties. They say little research has been done on the language. Anyone working on this language should use the dossier for ypg (Phola) and the entries for ypa and ypo.

\section{Primary Educational Resources}
None found. The language is not known to be taught in schools.

\section{Linguistic and Reference Materials}
\begin{itemize}
\item Glottolog entry for the retired code ypl: \url{https://glottolog.org/resource/languoid/iso/ypl}
\item Pelkey, ``Puzzling over Phula'' (academic survey): \url{https://www.academia.edu/2435607/Puzzling_over_Phula}
\item Phula numerals page, Max Planck Institute for Evolutionary Anthropology (linked to Ethnologue; I did not check its contents): \url{https://lingweb.eva.mpg.de/channumerals/Phula.htm}
\end{itemize}

\section{Religious and Community Resources}
\begin{itemize}
\item Asia Harvest profile of the Pula (an older ethnic-group profile, PDF): \url{https://people-groups.asiaharvest.org/China/chinaPeoples/P/Pula.pdf}
\item Joshua Project profile of Phola in China: \url{https://joshuaproject.net/index.php/people_groups/16151}
\end{itemize}

\section{Digital and Audio Resources}
None found for ypl.

\section{Curriculum Development Context}
Because the code is retired, any curriculum work belongs under Phola (ypg) and its related varieties Phala (ypa) and Alo Phola (ypo). No teaching materials were found.
\end{document}